% SPDX-FileCopyrightText: 2026 Yiyu Zhou <yiyu@yiyuzhou.io>
%
% SPDX-License-Identifier: GPL-3.0-or-later

\chapter{Lights}
\label{ch:lights}

Lighting consoles, dimmers, and nodes take levels over a network as sACN,
the open standard ANSI E1.31, or as Art-Net.  cinq sends both, so light cues
can set the lights along with the sound.  This chapter needs no lighting
hardware, since the window shows every level that cinq sends.

Press \key{Control+l} and open \texttt{lights.cinq.toml} in your copy of the
demo.

\section{The light patch}

Each light is an instrument with a name and a DMX address.  The address is
the channel that sets the intensity of the instrument.  The instruments form
the light patch.

\begin{enumerate}
\item Click the \ui{Light Patch} tab of the inspector.  It lists the three
  instruments of the show, as in Figure~\ref{fig:patch}.
\item Press \ui{GO} twice.  Cue~1 brings the front light to 30\% at once.
  Cue~2 fades it to 70\% and the back light to 40\% over three seconds.  The
  \ui{Now} column follows each level.
\end{enumerate}

\begin{figure}[H]
  \centering
  \includegraphics[width=\textwidth]{screenshots/patch.png}
  \caption{\ui{Light Patch} tab after cue~2}
  \label{fig:patch}
\end{figure}

Universes count from 1, as in sACN, and each holds 512 addresses.  The
\ui{Address} column shows an address in universe~1 as the address alone.  In
another universe, it shows the universe, a slash, and the address, as
\texttt{2/1} for the lightning.  \ui{Add instrument} adds an instrument at
the address after the last one.  \ui{Remove} removes an instrument from the
patch, and no cue sets it afterward.

\section{Light cues}

A light cue fades the instruments that it names to its own levels, in
percent, over its duration.  The other instruments keep their levels until a
later light cue names them, as in QLab.

\begin{enumerate}
\item Select cue~2, Stage up, and click the \ui{Levels} tab.  It lists every
  instrument with its level in the cue, as in Figure~\ref{fig:light}.  An
  empty field, as for the lightning, leaves the instrument unchanged.
\item Press \ui{GO}\@.  Cue~3 flashes the lightning.  It auto-continues to
  cue~4, which fades the lightning out in half a second, and cue~4
  auto-continues at once to cue~5, which plays the thunder.
\end{enumerate}

\begin{figure}[H]
  \centering
  \includegraphics[width=\textwidth]{screenshots/light.png}
  \caption{\ui{Levels} tab of a light cue}
  \label{fig:light}
\end{figure}

Pausing a light cue pauses its fade.  Stopping it, as Panic does, leaves the
lights where the fade has brought them, so Panic never blacks out the stage.
For a blackout, use a light cue that takes every instrument to 0\%, such as
cue~7.

\section{Seeing the output}

cinq sends each universe of the patch as sACN and as Art-Net.  It sends them
on every network of the computer and to the computer itself, whenever a level
changes and again every second.  An sACN monitor, such as sACNView, shows
universes 1 and 2 as the cues run.  cinq also announces its universes every
ten seconds, so the monitor lists the source \texttt{cinq} without further
setup.  On Darwin, allow cinq to find devices on local networks when it asks.

A receiver merges the sACN of cinq with that of any other source at the same
priority that drives the same universe, such as a lighting console, and takes
the higher level of each address.  Test on a universe that nothing else uses.
Drive the universe of a real rig only with the agreement of the lighting
designer.

\section*{Exercise}

Add an instrument with \ui{Add instrument}.  Then choose \ui{Light} in the
\ui{New cue} list and press \ui{Add cue}.  Set the \ui{Duration} of the new
cue to 2 seconds, and on the \ui{Levels} tab of the cue, set the level of the
new instrument to 50.  Press \ui{Standby selected}, then \ui{GO}\@.  Watch
the level in the \ui{Now} column or in your sACN monitor.
